\section{Relational witness clouds, backward viability, and stable erasure}
\label{sec:clouds}

Exact witness choice is one of the recurring obstructions in the repository.  Polynomial roots collide,
eigenvectors are nonunique, local analytic branches merge, minimizers need not be isolated, and finite
covers may admit many compatible refinements.  The right internal object is often a finite relation of
certified possibilities rather than a purported canonical selector.

\subsection{Layered finite clouds}

Fix a finite horizon \(h\).  At level \(i\), let candidates be codes
\[
 C_i=\{c<B_i : V_i(c)=1\},
\]
where \(B_i\) and \(V_i\) are primitive recursive.  Adjacent levels have a decidable compatibility relation
\(E_i(c,d)\).  A valid path is a sequence
\[
 c_0\in C_0,\ldots,c_h\in C_h,
 \qquad E_i(c_i,c_{i+1})	ext{ for }i<h.
\]
Nonemptiness of every \(C_i\) is not enough.  The candidates may be disconnected across levels.  This
simple observation invalidates several informal ``choose compatible approximations'' arguments unless a
compatibility theorem is supplied.

\subsection{Backward viability}

Define viable sets recursively from the final level:
\begin{align*}
 W_h(c) &:= V_h(c),\\
 W_i(c) &:= V_i(c)\wedge
    \exists d<B_{i+1}\bigl(E_i(c,d)\wedge W_{i+1}(d)\bigr).
\end{align*}
All quantifiers are bounded, so Theorem~\ref{thm:booleanization} produces concrete Boolean checkers.  The
tables can be computed by reverse primitive recursion over the fixed horizon.  A candidate is viable exactly
when it extends to a complete valid path.

\begin{theorem}[Finite viability]
\label{thm:viability}
For each \(i\le h\) and candidate \(c<B_i\), \(W_i(c)=1\) if and only if there exists a valid path from
\(c\) at level \(i\) to level \(h\).  If some \(c_0\) satisfies \(W_0(c_0)=1\), bounded least search
constructs a complete valid path.
\end{theorem}

\begin{proof}
Use reverse induction on \(i\).  At \(h\), a path of length zero is exactly a valid final candidate.  For the
step, \(W_i(c)=1\) says that \(c\) is valid and has a compatible viable successor \(d\).  The reverse
induction hypothesis turns viability of \(d\) into a path from \(d\) to the final level; prepend \(c\).
Conversely, any path from \(c\) has a first successor \(d\), and the tail path implies \(W_{i+1}(d)=1\), so
\(W_i(c)=1\).  Once \(W_0\) is nonempty, choose the least viable initial code.  At each next level choose
the least compatible viable successor by bounded search.  The defining equation of \(W_i\) guarantees that
the search succeeds.  Iterating for \(h\) steps constructs the path, and the invariant records validity and
compatibility of the prefix.
\end{proof}

The theorem is entirely finitary.  It uses no dependent choice, no compactness, and no infinite path
principle.  When an application needs an infinite coherent name, it must provide a quantitative method for
turning a sequence of finite-horizon paths into finite observations; this is addressed below.

\subsection{Exact branch-independent observables}

Suppose each valid final path \(p\) has an observation \(O(p)\), and the object theory proves
\[
  \operatorname{Path}(p)\wedge\operatorname{Path}(p')	o O(p)=O(p').
\]
Then choose the least viable path and return its observation.  The path is an internal implementation
detail.  The independence theorem proves that every other valid path gives the same public value.

\begin{theorem}[Exact branch erasure]
A nonempty finite path relation plus object-theory branch independence yields a deterministic public
observable.
\end{theorem}
\begin{proof}
The viability theorem supplies one path \(p_0\).  Define the public producer to return \(O(p_0)\).  For an
arbitrary valid path \(p\), branch independence yields \(O(p_0)=O(p)\).  Hence the public theorem is
independent of the least-code convention.  At proof-term level, instantiate the comparison theorem at
\(p_0,p\) and cut with their validity proofs.
\end{proof}

Typical exact stable observables include a polynomial's elementary symmetric functions, the dimension of a
finite vector space, a determinant, a multiset of roots, a homology rank, or a quotient object represented
by any of several equivalent finite presentations.

\subsection{Diameter-stable observables}

Exact independence is often too strong.  Let observations be rational codes with a primitive-recursive
metric \(d\).  Suppose the cloud has certified diameter
\[
  \operatorname{Path}(p)\wedge\operatorname{Path}(p')	o
  d(O(p),O(p'))\le\varepsilon(q).
\]
Selecting the least viable path gives an approximation determined up to \(\varepsilon(q)\).  This is enough
for represented real or complex outputs, approximate spectra, root clusters, and finite-net minimizers.

\begin{theorem}[Diameter erasure]
If every valid path observation lies within rational diameter \(\varepsilon(q)\), then any bounded path
selector is a valid \(\varepsilon(q)\)-observation of the cloud.  Two independent implementations of the
selector return outputs within \(\varepsilon(q)\).
\end{theorem}
\begin{proof}
Both selected paths satisfy the path verifier.  Apply the diameter theorem.  No equality of paths or
coherent choice is required.
\end{proof}

The public formula should state the metric estimate, not exact equality.  Converting a diameter theorem into
an exact label is valid only when an additional discrete separation bound is available.

\subsection{Exact recovery under a gap}

Assume candidate labels lie in a discrete finite set and distinct labels have observable separation greater
than \(2\varepsilon\).  If the cloud diameter is at most \(\varepsilon\), all candidates must have the same
label.  More generally, suppose a code \(\ell\) has a certified gap
\[
  d(z,Z_\ell)\le\varepsilon,
  \qquad
  d(Z_\ell,Z_j)>2\varepsilon\quad(j\ne\ell).
\]
Then every valid observation belongs to label \(\ell\).

\begin{proposition}[Gap-based branch collapse]
Approximation data plus an explicit separation certificate yields exact branch identification.
\end{proposition}
\begin{proof}
If another label \(j\ne\ell\) were valid, choose observations \(z_\ell,z_j\) for the two labels.  The
triangle inequality and the approximation bounds would give
\(d(Z_\ell,Z_j)\le 2\varepsilon\), contradicting the gap.  All quantities are rationally coded, so the
contradiction is a finite integer inequality after clearing denominators.
\end{proof}

This proposition is the legitimate form of ``precision eventually determines the branch.''  Without a
separation certificate the conclusion is false near collisions.

\subsection{Shrinking clouds and coherent names}

For each precision \(q\), suppose there is a nonempty finite cloud \(C_q\) and a rational observation
\(a_q(c)\).  There are two sufficient routes to a represented object.

\paragraph{Common-center diameter.}
If the object theory proves
\[
 \diam\{a_q(c):c\in C_q\}\le 2^{-q}
\]
and a cross-level compatibility estimate
\[
 c\in C_q,\quad d\in C_{q+1}
 \quad\Longrightarrow\quad
 |a_q(c)-a_{q+1}(d)|\le K2^{-q}.
\]
displayed for a fixed rational \(K\), then least-code selections form a Cauchy name with an explicit
modulus.

\paragraph{Summable motion.}
Alternatively, finite viability may produce a path with
\[
 |a_{q+1}(c_{q+1})-a_q(c_q)|\le 2^{-q-2}.
\]
Then for \(r>q\), a finite geometric-sum calculation gives
\[
 |a_r-a_q|\le \sum_{i=q}^{r-1}2^{-i-2}<2^{-q-1}.
\]
The relevant theorem is a finite inequality for arbitrary \(q,r\), proved by induction; no completed
infinite sum is required in iPRA.

In both routes the selector is primitive recursive because each level is bounded.  The public meaning is
determined by the diameter or motion theorem, not by the arbitrary least-code rule.

\subsection{Root clusters}

Let \(p(X)\) be a polynomial with exact rational or algebraic coefficients.  Near a multiple root, ordering
individual roots continuously is impossible.  A finite root-cluster contract instead returns disjoint
rational rectangles or disks, each with a multiplicity certificate and a completeness theorem.  As
precision increases, clusters may split but must refine compatibly.  Stable observations include:
\begin{itemize}
  \item the multiset of clusters with multiplicities;
  \item symmetric polynomials of all roots in a cluster;
  \item the total number of roots in a region;
  \item an approximate Hausdorff set of roots;
  \item an individually labeled root only after a separation gap is certified.
\end{itemize}
Backward viability handles the refinement relation.  This directly resolves the exact-root-selection
problem noted in Puiseux, spectral, and polynomial-factorization packets: do not demand a global label where
mathematics supplies only a finite relation.

\subsection{Finite covers and gluing}

A cover problem fits the same pattern.  Level \(i\) may list local sections or chart data at precision or
refinement \(i\); compatibility edges assert equality on overlaps or controlled discrepancy.  Backward
viability finds a finite compatible family.  A gluing theorem then maps any compatible family to a global
finite observation.  If gluing is unique, the result is exact branch independent.  If only approximate
uniqueness holds, the output is diameter stable.

The framework does not infer a finite subcover from arbitrary compactness.  The cover list or a bound on it
must be supplied.  Once supplied, all compatibility and gluing checks are finite and primitive recursive.

\subsection{What finite viability does not prove}

It is important not to overread Theorem~\ref{thm:viability}.
\begin{itemize}
  \item Viability for every finite horizon, with no uniform bound or compatibility modulus, does not by
        itself produce one infinite path in iPRA.
  \item Nonempty levels do not imply a path.
  \item Small cardinality does not imply that candidate codes have a bounded height.
  \item A least-code path need not be mathematically continuous in parameters.
\end{itemize}
What is available is exactly what finite consumers need: for any supplied finite horizon, a verified path
and a stable observation.  When a full represented object is required, a shrinking-diameter or summable-
motion theorem must be included.
